% ============================================================
\chapter{Étude et réalisation du Sprint 2}
\label{chap:sprint2}
% ============================================================

\section{Introduction}
\label{sec:intro_sprint2}

Ce chapitre présente l'étude et la réalisation du deuxième sprint de la
plateforme bancaire FinBank. Ce sprint constitue la base sécuritaire de
la solution. Il concerne la gestion des utilisateurs, l'authentification,
les rôles, la gestion du profil et la protection des accès.

La première partie présente le backlog du sprint. La deuxième partie
décrit la conception à travers les diagrammes UML et les descriptions
textuelles des principaux cas d'utilisation. La troisième partie
présente les interfaces réalisées et explique le rôle de chaque capture
insérée.

\section{Backlog du Sprint}
\label{sec:backlog_sprint2}

Le tableau~\ref{tab:backlog_sprint2} présente les tâches réalisées
durant le Sprint~2 ainsi que leur estimation.

\begin{table}[htbp]
\centering
\caption{Backlog du Sprint 2}
\label{tab:backlog_sprint2}
\renewcommand{\arraystretch}{1.35}
\begin{tabular}{|c|p{9.5cm}|c|}
\hline
\textbf{ID} & \textbf{Tâche} & \textbf{Estimation} \\
\hline
S2-T1 & Mise en place du backend et de la configuration de la persistance & 2 jours \\
\hline
S2-T2 & Développement de l'inscription, de la connexion, du renouvellement de session et de la déconnexion & 4 jours \\
\hline
S2-T3 & Gestion des utilisateurs et des rôles applicatifs (\texttt{USER}, \texttt{ADMIN}) & 3 jours \\
\hline
S2-T4 & Protection des ressources selon le rôle et l'état du compte & 3 jours \\
\hline
S2-T5 & Gestion du profil et changement du mot de passe & 2 jours \\
\hline
S2-T6 & Activation et vérification de la double authentification (2FA/TOTP) & 2 jours \\
\hline
S2-T7 & Réalisation des interfaces d'inscription, de connexion et de profil & 3 jours \\
\hline
S2-T8 & Vérification des fonctionnalités de sécurité et correction des anomalies & 2 jours \\
\hline
\multicolumn{2}{|r|}{\textbf{Total}} & \textbf{21 jours} \\
\hline
\end{tabular}
\end{table}

\section{Conception}
\label{sec:conception_sprint2}

La conception du Sprint~2 permet de représenter les acteurs, les
fonctionnalités et les classes nécessaires à la sécurité de FinBank.
Elle est organisée autour du diagramme de cas d'utilisation, du
diagramme de classes, des descriptions textuelles, des classes
participantes et des diagrammes de séquence.

\subsection{Diagramme de cas d'utilisation}
\label{subsec:usecase_sprint2}

Le diagramme de la figure~\ref{fig:usecase_sprint2} présente les
interactions entre les acteurs et le système pour ce sprint. Les
acteurs utilisés correspondent aux rôles réellement présents dans
l'application : l'\textbf{Utilisateur (USER)} et
l'\textbf{Administrateur (ADMIN)}.

\begin{figure}[htbp]
\centering
\IfFileExists{Figures/usecase_sprint2.png}{%
  \includegraphics[width=\textwidth,height=0.76\textheight,keepaspectratio]{Figures/usecase_sprint2.png}%
}{%
  \fbox{\parbox[c][5cm][c]{0.88\textwidth}{\centering
  Espace réservé au diagramme de cas d'utilisation du Sprint 2.}}
}
\caption{Diagramme de cas d'utilisation du Sprint 2}
\label{fig:usecase_sprint2}
\end{figure}

\noindent
L'utilisateur peut créer un compte, s'authentifier, gérer son profil et
activer la double authentification. Après une connexion réussie, il
peut accéder uniquement aux fonctionnalités autorisées par son rôle.
L'administrateur dispose en plus des fonctions de gestion des
utilisateurs. Les fonctionnalités protégées nécessitent une
authentification préalable.

\subsection{Diagramme de classes}
\label{subsec:classes_sprint2}

Le diagramme de classes de la figure~\ref{fig:classes_sprint2}
présente les principales entités du périmètre du Sprint~2. Il décrit
les informations d'identité, les rôles, les états du compte et les
éléments utilisés pour gérer les sessions et la récupération du mot de
passe.

\begin{figure}[htbp]
\centering
\IfFileExists{Figures/classes_sprint2.png}{%
  \includegraphics[width=\textwidth,height=0.78\textheight,keepaspectratio]{Figures/classes_sprint2.png}%
}{%
  \fbox{\parbox[c][5cm][c]{0.88\textwidth}{\centering
  Espace réservé au diagramme de classes du Sprint 2.}}
}
\caption{Diagramme de classes du Sprint 2}
\label{fig:classes_sprint2}
\end{figure}

\noindent
La classe \texttt{User} occupe une position centrale dans le modèle.
Elle est associée aux rôles et aux états du compte. Les classes
\texttt{RefreshToken}, \texttt{PasswordResetToken} et
\texttt{AuthToken} permettent respectivement de gérer le renouvellement
de session, la récupération du mot de passe et la persistance des
informations d'authentification. Cette modélisation sépare les données
de l'utilisateur des mécanismes de sécurité qui lui sont associés.

\subsection{Description textuelle des cas d'utilisation}
\label{subsec:desc_cas_sprint2}

Les tableaux suivants précisent les acteurs, les préconditions, le
scénario nominal et les scénarios alternatifs des deux cas
d'utilisation principaux du sprint.

\subsubsection{S'authentifier}
\label{subsubsec:desc_auth_sprint2}

\begin{table}[htbp]
\centering
\caption{Description textuelle -- S'authentifier}
\label{tab:uc_auth_sprint2}
\renewcommand{\arraystretch}{1.3}
\begin{tabular}{|p{3.1cm}|p{10.6cm}|}
\hline
\textbf{Élément} & \textbf{Description} \\
\hline
Acteur & Utilisateur (\texttt{USER}) ou Administrateur (\texttt{ADMIN}) \\
\hline
But & Permettre l'accès à l'espace correspondant au rôle de l'utilisateur. \\
\hline
Précondition & L'utilisateur possède un compte valide et actif. \\
\hline
Postcondition & L'utilisateur est authentifié et peut accéder aux fonctionnalités autorisées. \\
\hline
Scénario principal &
1. L'utilisateur ouvre l'écran de connexion. \newline
2. Il saisit son adresse e-mail et son mot de passe. \newline
3. Le système vérifie les informations. \newline
4. Le système identifie le rôle et l'état du compte. \newline
5. L'accès est accordé. \\
\hline
Scénarios alternatifs & Si les informations sont incorrectes, un message d'erreur est affiché. Si la double authentification est activée, le système demande un code supplémentaire. \\
\hline
\end{tabular}
\end{table}

\subsubsection{Gérer son profil}
\label{subsubsec:desc_profil_sprint2}

\begin{table}[htbp]
\centering
\caption{Description textuelle -- Gérer son profil}
\label{tab:uc_profil_sprint2}
\renewcommand{\arraystretch}{1.3}
\begin{tabular}{|p{3.1cm}|p{10.6cm}|}
\hline
\textbf{Élément} & \textbf{Description} \\
\hline
Acteur & Utilisateur (\texttt{USER}) ou Administrateur (\texttt{ADMIN}) \\
\hline
But & Consulter et modifier les informations personnelles du compte. \\
\hline
Précondition & L'utilisateur est authentifié. \\
\hline
Postcondition & Les informations valides sont enregistrées. \\
\hline
Scénario principal &
1. L'utilisateur ouvre son profil. \newline
2. Le système affiche les informations disponibles. \newline
3. L'utilisateur modifie les champs souhaités. \newline
4. Le système vérifie les données. \newline
5. Les modifications sont enregistrées et confirmées. \\
\hline
Scénario alternatif & Si une donnée est absente ou invalide, le système refuse la modification et affiche une indication corrective. \\
\hline
\end{tabular}
\end{table}

\subsection{Diagrammes de classes participantes}
\label{subsec:participants_sprint2}

Les diagrammes de classes participantes détaillent les collaborations
nécessaires à l'exécution des cas d'utilisation. Ils complètent le
diagramme de classes en montrant les rôles de l'interface, du contrôle,
des services et des entités.

\subsubsection{Classes participantes de l'authentification}
\label{subsubsec:participants_auth_sprint2}

\begin{figure}[htbp]
\centering
\IfFileExists{Figures/participants_auth.png}{%
  \includegraphics[width=\textwidth,height=0.74\textheight,keepaspectratio]{Figures/participants_auth.png}%
}{%
  \fbox{\parbox[c][5cm][c]{0.88\textwidth}{\centering
  Espace réservé au diagramme des classes participantes de l'authentification.}}
}
\caption{Classes participantes -- Authentification}
\label{fig:participants_auth}
\end{figure}

\noindent
La figure~\ref{fig:participants_auth} montre la collaboration entre
l'utilisateur, l'interface de connexion, le composant de contrôle, le
service d'authentification et les éléments persistants. Cette
organisation permet de séparer la saisie des informations, les règles
de sécurité et l'accès aux données.

\subsubsection{Classes participantes de la gestion du profil}
\label{subsubsec:participants_profil_sprint2}

\begin{figure}[htbp]
\centering
\IfFileExists{Figures/participants_profil.png}{%
  \includegraphics[width=\textwidth,height=0.74\textheight,keepaspectratio]{Figures/participants_profil.png}%
}{%
  \fbox{\parbox[c][5cm][c]{0.88\textwidth}{\centering
  Espace réservé au diagramme des classes participantes du profil.}}
}
\caption{Classes participantes -- Gestion du profil}
\label{fig:participants_profil}
\end{figure}

\noindent
La figure~\ref{fig:participants_profil} présente le cheminement d'une
modification du profil. L'interface transmet les informations, le
contrôle vérifie l'utilisateur authentifié, le service applique les
règles métier et la couche de persistance enregistre uniquement les
données valides.

\subsection{Diagrammes de séquence}
\label{subsec:sequences_sprint2}

Les diagrammes de séquence décrivent le déroulement temporel de
l'authentification selon l'état de la double authentification.

\subsubsection{Authentification sans double authentification}
\label{subsubsec:seq_sans_2fa_sprint2}

\begin{figure}[htbp]
\centering
\IfFileExists{Figures/seq_auth_sans_2fa.png}{%
  \includegraphics[width=\textwidth,height=0.76\textheight,keepaspectratio]{Figures/seq_auth_sans_2fa.png}%
}{%
  \fbox{\parbox[c][5cm][c]{0.88\textwidth}{\centering
  Espace réservé au diagramme de séquence sans double authentification.}}
}
\caption{Diagramme de séquence -- Authentification sans 2FA}
\label{fig:seq_auth_sans_2fa}
\end{figure}

\noindent
Dans ce scénario, l'utilisateur fournit ses identifiants. Le système
recherche le compte, vérifie le mot de passe et contrôle l'état du
compte. Lorsque les informations sont correctes, l'accès est accordé et
la session authentifiée est créée. Dans le cas contraire, le système
retourne une erreur sans ouvrir la session.

\subsubsection{Authentification avec double authentification}
\label{subsubsec:seq_avec_2fa_sprint2}

\begin{figure}[htbp]
\centering
\IfFileExists{Figures/seq_auth_avec_2fa.png}{%
  \includegraphics[width=\textwidth,height=0.78\textheight,keepaspectratio]{Figures/seq_auth_avec_2fa.png}%
}{%
  \fbox{\parbox[c][5cm][c]{0.88\textwidth}{\centering
  Espace réservé au diagramme de séquence avec double authentification.}}
}
\caption{Diagramme de séquence -- Authentification avec 2FA}
\label{fig:seq_auth_avec_2fa}
\end{figure}

\noindent
Lorsque la double authentification est activée, le traitement se
déroule en deux phases. Après la vérification des identifiants, le
système demande un code temporaire généré par l'application
d'authentification. L'accès définitif est accordé uniquement si ce code
est valide. Cette étape renforce la protection du compte même si le
mot de passe est connu.

\section{Réalisation}
\label{sec:realisation_sprint2}

La réalisation du Sprint~2 concerne les interfaces nécessaires à
l'inscription, à la connexion, à la double authentification et à la
gestion du profil. Chaque capture est accompagnée d'une explication
de son contenu et de sa contribution au fonctionnement de la
plateforme.

\subsection{Interface d'inscription}
\label{subsec:ui_register_sprint2}

\begin{figure}[htbp]
\centering
\IfFileExists{Figures/ui_register.png}{%
  \includegraphics[width=0.90\textwidth,keepaspectratio]{Figures/ui_register.png}%
}{%
  \fbox{\parbox[c][4cm][c]{0.85\textwidth}{\centering
  Espace réservé à la capture de l'interface d'inscription.}}
}
\caption{Interface d'inscription}
\label{fig:ui_register_sprint2}
\end{figure}

\noindent
Cette interface permet à un nouvel utilisateur de créer son compte.
Elle regroupe les informations nécessaires à l'inscription et fournit
des indications lorsque les données saisies ne respectent pas les
règles attendues. Après validation, le compte peut être utilisé pour
accéder à l'espace sécurisé de la plateforme.

\subsection{Interface de connexion}
\label{subsec:ui_login_sprint2}

\begin{figure}[htbp]
\centering
\IfFileExists{Figures/ui_login.png}{%
  \includegraphics[width=0.90\textwidth,keepaspectratio]{Figures/ui_login.png}%
}{%
  \fbox{\parbox[c][4cm][c]{0.85\textwidth}{\centering
  Espace réservé à la capture de l'interface de connexion.}}
}
\caption{Interface de connexion}
\label{fig:ui_login_sprint2}
\end{figure}

\noindent
Cette capture présente le point d'entrée de la plateforme. L'utilisateur
saisit ses informations d'accès, puis le système vérifie son identité
et l'oriente vers l'espace correspondant à son rôle. En cas d'activation
de la double authentification, l'étape de vérification supplémentaire
est affichée avant l'accès final.

\subsection{Interface de vérification 2FA}
\label{subsec:ui_2fa_sprint2}

\begin{figure}[htbp]
\centering
\IfFileExists{Figures/ui_2fa.png}{%
  \includegraphics[width=0.90\textwidth,keepaspectratio]{Figures/ui_2fa.png}%
}{%
  \fbox{\parbox[c][4cm][c]{0.85\textwidth}{\centering
  Espace réservé à la capture de l'interface de vérification 2FA.}}
}
\caption{Interface de vérification de la double authentification}
\label{fig:ui_2fa_sprint2}
\end{figure}

\noindent
Cette interface demande le code de vérification généré par
l'application d'authentification. Le code est contrôlé avant
l'ouverture de la session. Un code incorrect entraîne l'affichage d'un
message d'erreur et empêche l'accès à l'espace sécurisé.

\subsection{Interface de gestion du profil}
\label{subsec:ui_profile_sprint2}

\begin{figure}[htbp]
\centering
\IfFileExists{Figures/ui_profile.png}{%
  \includegraphics[width=0.90\textwidth,keepaspectratio]{Figures/ui_profile.png}%
}{%
  \fbox{\parbox[c][4cm][c]{0.85\textwidth}{\centering
  Espace réservé à la capture de l'interface du profil.}}
}
\caption{Interface de gestion du profil}
\label{fig:ui_profile_sprint2}
\end{figure}

\noindent
Cette interface permet à l'utilisateur authentifié de consulter et de
mettre à jour ses informations personnelles. Elle permet également
d'accéder aux paramètres de sécurité, notamment à la gestion du mot de
passe et de la double authentification. Les modifications sont
confirmées uniquement après validation des informations saisies.

\section{Conclusion}
\label{sec:conclusion_sprint2}

Ce chapitre a présenté l'étude et la réalisation du Sprint~2. Le
backlog a d'abord défini les tâches liées à la sécurité et à la gestion
des utilisateurs. La conception a ensuite été détaillée à travers les
diagrammes de cas d'utilisation, de classes, de classes participantes
et de séquence.

La réalisation a enfin présenté les interfaces d'inscription, de
connexion, de double authentification et de gestion du profil. Ce
sprint fournit ainsi une base sécurisée pour les fonctionnalités
bancaires développées dans le Sprint~3.
